% 自编教学补充；阅读来源见本章 README。
\begin{frame}[fragile]{1. 校园饮品订单：一行一笔订单}
\small 六笔合成订单用于手工核对。金额为数量乘单价；本例价格均已匹配，日期是实际有订单的日期。
\medskip
\begin{lstlisting}
import pandas as pd
sales = pd.DataFrame({
    "订单": [1, 2, 3, 4, 5, 6],
    "日期": ["2026-09-01"] * 3 + ["2026-09-02"] * 3,
    "门店": ["东", "东", "西", "东", "西", "西"],
    "商品": ["茶", "咖啡", "茶", "茶", "咖啡", "茶"],
    "数量": [2, 1, 1, 3, 2, 2], "单价": [10, 20, 10, 10, 20, 10]})
sales["日期"] = pd.to_datetime(sales["日期"])
sales["金额"] = sales["数量"] * sales["单价"]
sales
\end{lstlisting}
\end{frame}

\begin{frame}[fragile]{2. groupby：拆分、计算、合并}
\small 先按门店拆成小表，在各小表内求和，再组成汇总表。as\_index=False 使门店保留为普通列，便于继续合并。
\medskip
\begin{lstlisting}
store_total = sales.groupby("门店", as_index=False)["金额"].sum()
print(store_total)
print(sales.groupby(["门店", "商品"])["金额"].sum())
\end{lstlisting}
\end{frame}

\begin{frame}[fragile]{3. 命名聚合：明确每个指标的含义}
\small 总金额、订单数和平均每单金额是不同指标。命名聚合直接给结果列起名；金额均值的分母是订单行数。
\medskip
\begin{lstlisting}
summary = sales.groupby("门店", as_index=False).agg(
    总金额=("金额", "sum"),
    订单数=("订单", "size"),
    总杯数=("数量", "sum"),
    平均每单金额=("金额", "mean"))
print(summary)
\end{lstlisting}
\end{frame}

\begin{frame}[fragile]{4. size 与 count：分母不能混淆}
\small size 统计组内行数，count 统计某列非缺失值数。先明确要数订单、有效金额还是不同顾客，再选择统计方法。
\medskip
\begin{lstlisting}
with_missing = sales.copy()
with_missing.loc[0, "金额"] = float("nan")
counts = with_missing.groupby("门店").agg(
    行数=("金额", "size"), 有效金额数=("金额", "count"))
print(counts)
\end{lstlisting}
\end{frame}
\begin{frame}[fragile]{课堂练习：4. size 与 count：分母不能混淆}
东店的行数与有效金额数各是多少？金额缺失能否当作免费订单？
\medskip
\begin{block}{检查思路}
先写出预期结果，再运行。参考答案见配套 Notebook 的折叠单元格。
\end{block}
\end{frame}

\begin{frame}[fragile]{5. 加权均值：平均每杯价格}
\small 平均每单金额不等于平均每杯价格。后者用总金额除以总杯数，相当于按购买数量给单价加权。
\medskip
\begin{lstlisting}
summary["平均每杯价格"] = summary["总金额"] / summary["总杯数"]
print(summary[["门店", "平均每单金额", "平均每杯价格"]])
print("全体每杯均价：", sales["金额"].sum() / sales["数量"].sum())
\end{lstlisting}
\end{frame}

\begin{frame}[fragile]{6. transform：把组内统计带回原行}
\small agg 每组输出一行，transform 返回与原表等长且对齐的结果。用每笔金额除以所属门店总金额得到组内占比。
\medskip
\begin{lstlisting}
sales["门店总金额"] = sales.groupby("门店")["金额"].transform("sum")
sales["门店内占比"] = sales["金额"] / sales["门店总金额"]
print(sales[["订单", "门店", "金额", "门店内占比"]])
print(sales.groupby("门店")["门店内占比"].sum())
\end{lstlisting}
\end{frame}

\begin{frame}[fragile]{7. 排名与组内前两名}
\small 先按金额降序，再按订单号升序决定并列次序，组内取前两行。结果是两笔订单，而不一定是两个不同金额档位。
\medskip
\begin{lstlisting}
ordered = sales.sort_values(["门店", "金额", "订单"],
                             ascending=[True, False, True])
top2 = ordered.groupby("门店").head(2)
print(top2[["订单", "门店", "金额"]])
sales["组内名次"] = sales.groupby("门店")["金额"].rank(
    ascending=False, method="min")
\end{lstlisting}
\end{frame}

\begin{frame}[fragile]{8. 交叉表：计数与金额表不同}
\small crosstab 默认数记录。pivot\_table 显式指定金额与 sum 才是金额汇总。本例已知交易明细完整，因此没有订单的组合可填零。
\medskip
\begin{lstlisting}
print(pd.crosstab(sales["门店"], sales["商品"]))
amount_table = sales.pivot_table(index="门店", columns="商品",
    values="金额", aggfunc="sum", fill_value=0)
print(amount_table)
\end{lstlisting}
\end{frame}

\begin{frame}[fragile]{9. 按日汇总与滚动窗口}
\small 先汇总到“日”，再滚动才是相邻观测日的窗口。rolling(2) 是两行，不自动补齐日历中没有记录的日期。
\medskip
\begin{lstlisting}
daily = sales.groupby("日期")["金额"].sum().sort_index()
print(daily)
print(daily.rolling(2, min_periods=2).mean())
print(sales.set_index("日期")["金额"].resample("D").sum(min_count=1))
\end{lstlisting}
\end{frame}

\begin{frame}[fragile]{10. 金融迁移：按股票计算收益率}
\small 合成长表每行一只股票一天的收盘价。先按股票和日期排序，在每只股票内部错位，避免把另一只股票的价格当作昨日价。
\medskip
\begin{lstlisting}
quotes = pd.DataFrame({
    "股票": ["A", "B", "A", "B", "A", "B"],
    "日期": pd.to_datetime(["2026-09-01"] * 2 +
                          ["2026-09-02"] * 2 + ["2026-09-03"] * 2),
    "收盘价": [100, 50, 102, 49, 101, 50]})
quotes = quotes.sort_values(["股票", "日期"])
quotes["昨日价"] = quotes.groupby("股票")["收盘价"].shift(1)
quotes["收益率"] = quotes["收盘价"] / quotes["昨日价"] - 1
print(quotes)
\end{lstlisting}
\end{frame}

\begin{frame}[fragile]{11. 金融迁移：从长表到收益率矩阵}
\small pivot 后每列一只股票；pct\_change(fill\_method=None) 不自动填补缺失价格。比较两种写法的同一观测值。
\medskip
\begin{lstlisting}
price_matrix = quotes.pivot(index="日期", columns="股票", values="收盘价")
return_matrix = price_matrix.pct_change(fill_method=None)
print(return_matrix)
expected = quotes.loc[(quotes["股票"] == "A") &
    (quotes["日期"] == "2026-09-02"), "收益率"].iloc[0]
assert abs(return_matrix.loc["2026-09-02", "A"] - expected) < 1e-12
\end{lstlisting}
\end{frame}

\begin{frame}[fragile]{12. 综合练习与验收}
\small 先用订单手算，再检查汇总表。全表金额应为 140 元、总杯数为 11；两门店各 70 元，但每杯均价不同。
\medskip
\begin{lstlisting}
assert sales["金额"].sum() == 140
assert sales["数量"].sum() == 11
assert summary["总金额"].sum() == sales["金额"].sum()
print(summary)
\end{lstlisting}
\end{frame}
\begin{frame}[fragile]{课堂练习：12. 综合练习与验收}
按“日期—门店”统计金额与杯数，并计算每杯均价。哪一天全体销售金额更高？
\medskip
\begin{block}{检查思路}
先写出预期结果，再运行。参考答案见配套 Notebook 的折叠单元格。
\end{block}
\end{frame}

\begin{frame}[fragile]{13. 综合案例起点：从原始订单到日报}
\small 下面重新开始一套独立的合成数据。每行原本应是一笔订单，但有重复登记、数量缺失和未知商品；目标是输出可核对的已知金额日报。
\medskip
\begin{lstlisting}
raw_orders = pd.DataFrame({
    "订单": ["01", "02", "03", "04", "05", "06", "01"],
    "日期": ["2026-09-01"] * 3 + ["2026-09-02"] * 3 + ["2026-09-01"],
    "门店": ["东", "东", "西", "东", "西", "西", "东"],
    "商品": [" t ", "C", "T", "T", "C", "X", " t "],
    "数量": ["2", "1", "缺失", "3", "2", "1", "2"]})
catalog = pd.DataFrame({"商品": ["T", "C"], "单价": [10, 20]})
raw_orders
\end{lstlisting}
\end{frame}

\begin{frame}[fragile]{14. 清洗：原始值与统计值并存}
\small 保留 raw\_orders。先统一商品代码，再删除完全相同的登记，转换日期和数量；原始数量列保留，便于回查缺失原因。
\medskip
\begin{lstlisting}
tidy_orders = raw_orders.copy()
tidy_orders["商品"] = tidy_orders["商品"].str.strip().str.upper()
tidy_orders = tidy_orders.drop_duplicates().reset_index(drop=True)
tidy_orders["日期"] = pd.to_datetime(tidy_orders["日期"])
tidy_orders["有效数量"] = pd.to_numeric(tidy_orders["数量"], errors="coerce")
assert tidy_orders["订单"].is_unique
print(tidy_orders)
\end{lstlisting}
\end{frame}

\begin{frame}[fragile]{15. 合并：把待核订单单独列出}
\small 左连接保留每笔订单。已匹配商品且数量为正才进入本例金额统计；数量缺失和价格未知都保留在待核清单中。
\medskip
\begin{lstlisting}
enriched = tidy_orders.merge(catalog, on="商品", how="left",
    validate="many_to_one", indicator=True)
eligible = enriched["_merge"].eq("both") & enriched["有效数量"].gt(0)
pending = enriched.loc[~eligible].copy()
known = enriched.loc[eligible].copy()
known["金额"] = known["有效数量"] * known["单价"]
print(pending[["订单", "数量", "商品", "_merge"]])
\end{lstlisting}
\end{frame}

\begin{frame}[fragile]{16. 汇总：金额与覆盖范围一起报告}
\small 已知金额日报仅覆盖四笔订单，不能称为全部销售额。汇总前后金额应一致，同时报告两笔待核订单。
\medskip
\begin{lstlisting}
known_report = known.groupby(["日期", "门店"], as_index=False).agg(
    已知金额=("金额", "sum"), 可计算订单数=("订单", "size"))
print(known_report)
print("待核订单数：", len(pending))
assert len(raw_orders) == 7 and len(tidy_orders) == 6
assert len(known) == 4 and len(pending) == 2
assert known_report["已知金额"].sum() == 110
\end{lstlisting}
\end{frame}

\begin{frame}[fragile]{17. 变形：缺失单元格不能随意填零}
\small 日报透视成“日期×门店”方便比较。西店首日有数量待核订单，不能把该格写成零销售；即使已有金额的格子也可能尚未完整。
\medskip
\begin{lstlisting}
known_wide = known_report.pivot(index="日期", columns="门店", values="已知金额")
pending_wide = pending.groupby(["日期", "门店"]).size().unstack("门店")
print("已知金额（不是完整销售额）：")
print(known_wide)
print("待核订单数：")
print(pending_wide)
\end{lstlisting}
\end{frame}
\begin{frame}[fragile]{课堂练习：17. 变形：缺失单元格不能随意填零}
核实订单03数量为1、商品X是5元的水。重新清洗合并并汇总，完整金额应为125元。不要在最终报表中直接改金额。
\medskip
\begin{block}{检查思路}
先写出预期结果，再运行。参考答案见配套 Notebook 的折叠单元格。
\end{block}
\end{frame}
